\section{Cache reuse and possible batching mechanisms}
\label{app:batching_architectural}

Two different forms of cache reuse matter: sharing unchanged context
across requests, and preserving activations within a request as generation
progresses. In AR decoding, causal attention allows the already-generated
prefix's KV state to persist while new tokens are appended. During
bidirectional denoising, changing canvas tokens can invalidate activations
within that canvas. This can increase the work per refinement step.

It does not follow that diffusion must discard all cached context.
DiffusionGemma's documented encoder--decoder design caches the prompt and
completed canvases, while the active canvas uses bidirectional attention
and accesses encoded context~\cite{diffusiongemma}. A changing canvas
also does not require a changing attention mask at every step.

The observations are therefore compatible with several explanations for
modest batching gains: larger per-request kernels, different weight and
activation reuse, scheduling restrictions, and implementation-specific
cache behavior. The existing controls do not distinguish their shares.
We do not predict a maximum batching multiplier or claim compute
saturation at batch size one. Establishing such a mechanism would require
hardware profiling and matched cache/scheduling ablations on a working
optimized diffusion path.
